\scanpage{264}
§ 2 Application à la géométrie algébrique sur une algèbre analytique.

\uline{Th}. 2.1. Soient $A$ une algèbre analytique sur $\mathbb C$, \inserted{$Y=\operatorname{Spec}A$,} $X$ un schéma propre sur $A$, \deleted{Alors} $\widehat Y=\operatorname{Specf}\widehat A$, $\widehat X=X\times_Y\widehat Y$, \deleted{Alors} de sorte que [cf] \deleted{$\unsure{\varphi}$} $\operatorname{Pic}(\widehat X)\simeq\varprojlim\operatorname{Pic}(X_n)$.

Alors

a) $(\operatorname{Pic}(X_n))_n$ satisfait ML

\margin{se généralise aux \uline{\unsure{schémas}} formels en \uuline{car. 0}
\figure{\begin{tikzpicture}[x=1cm,y=1cm]\draw[<->] (0,0)--(2.7,.7);\end{tikzpicture}}}

(ii) $\operatorname{Pic}(X)$ et $\operatorname{Pic}(\widehat X)$ ont $=$ image dans $\operatorname{Pic}(X_{\unsure{0|n}})$.
[\unsure{donc les comp. universelles|donne les comp. universelles}]

(iii) Pour q
\[
\varphi:\operatorname{Pic}(X)\unsure{-|\to}\operatorname{Pic}(\widehat X)
\]
soit un isom il faut s. que $H^1(X,\mathcal O_X)$ soit de longueur finie sur $A$, i.e. $\operatorname{supp}R^1f_*(\mathcal O_X)\subset\{y\}$ ($y$ le pt fermé de $Y$), \deleted{\unsure{ce qui équivaut|ce qui implique} \inserted{\unsure{le groupe}} \unsure{de} \unsure{isomorphisme} \unsure{que le syst.} projectif $(\operatorname{Pic}(X_n))$ soit ess\textsuperscript{t}. constant}
\figure{\begin{tikzpicture}[x=1cm,y=1cm]
\draw (0,0)--(2.5,.8)--(5.2,1.3)..controls(6,1.4)and(6.6,1.3)..(7,1.3)..controls(7.3,1.2)and(7.2,.8)..(7.8,.65)--(10,.7)--(10.5,.5)--(10.4,-.1)--(0,-.2)--(0,0);
\draw (0,.1)--(9.5,1.1);\draw (0,-.1)--(10.5,.5);
\draw (0,0)..controls(.8,-.4)and(1.1,.9)..(1.2,.2)..controls(1.5,-.2)and(1.8,.4)..(2,.9)..controls(2.2,-.4)and(2.5,-.1)..(2.7,.1)--(3.2,1)..controls(3.4,-.3)and(3.8,.1)..(4.1,.7)..controls(4.3,-.3)and(4.7,-.2)..(5,.3)--(5.2,1.2)..controls(5.3,-.2)and(5.8,.5)..(6,1.1)..controls(6.3,.3)and(6,.0)..(6.8,.1)--(7.5,.9)..controls(7.6,-.1)and(8.2,-.1)..(8.6,.0)--(8.5,.8)--(9.1,.2)--(9.8,-.1)--(10.2,.5);
\draw (5.3,1.2) ellipse (.9,.25);
\end{tikzpicture}}

Notons de suite le corollaire suivant de \deleted{l'injectivité de} a) et ii)

\uline{Corollaire} 2.2. Conditions équivalentes

\mspage{263}
